\documentclass[aspectratio=169, 14pt]{beamer}
\usepackage{stampfly_slides}

\lesson{3}{LED 制御}{LED Control}

\begin{document}

% ------------------------------------------------------------------
\begin{frame}
  \titlepage
\end{frame}

% ------------------------------------------------------------------
\begin{frame}{今日のゴール / Today's Goal}
  \begin{block}{目標}
    LED でシステムの状態（ARM / DISARM）とバッテリー残量を可視化する
  \end{block}

  \vspace{1em}

  \begin{itemize}
    \item WS2812 アドレサブル LED の制御
    \item 状態遷移の考え方（ステートマシン）
    \item バッテリー電圧のモニタリング
  \end{itemize}
\end{frame}

% ------------------------------------------------------------------
\begin{frame}{LED 状態遷移 / LED State Machine}
  \begin{center}
    \resizebox{0.95\linewidth}{!}{\documentclass[border=10pt]{standalone}
\usepackage{tikz}
\usetikzlibrary{arrows.meta, positioning, automata, calc}

\begin{document}

\definecolor{redstate}{RGB}{220, 50, 30}
\definecolor{greenstate}{RGB}{40, 160, 40}
\definecolor{gray1}{RGB}{120, 120, 120}
\begin{tikzpicture}[
    state/.style={
      circle, draw=#1, line width=2pt, fill=#1!10,
      minimum size=32mm, font=\bfseries\large, text=#1, align=center
    },
    trans/.style={-{Stealth[length=3.5mm]}, line width=1.5pt, gray1},
    tlbl/.style={font=\small, text=gray1, align=center, fill=white, inner sep=2pt},
    every initial by arrow/.style={trans},
  ]

  % States
  \node[state=redstate] (disarm) at (0, 0)
    {DISARM\\[2pt]\scriptsize\textcolor{gray1}{Battery gradient}\\
     \scriptsize\textcolor{redstate}{3.3\,V = Red}
     \textcolor{gray1}{$\rightarrow$}
     \textcolor{greenstate}{4.2\,V = Green}};
  \node[state=greenstate, right=55mm of disarm] (arm)
    {ARM\\\textcolor{greenstate!70}{\rule{8pt}{8pt}} Green};

  % Transitions
  \draw[trans] (disarm) to[bend left=20]
    node[tlbl, above] {is\_armed() == true} (arm);
  \draw[trans] (arm) to[bend left=20]
    node[tlbl, below] {is\_armed() == false} (disarm);

  % Initial
  \draw[trans] (-3, 0) -- (disarm)
    node[midway, above, font=\small, text=gray1] {Power ON};

  % LED color labels
  \node[font=\scriptsize, text=gray1, below=1mm of disarm]
    {led\_color(r, g, 0)};
  \node[font=\scriptsize, text=greenstate, below=1mm of arm]
    {led\_color(0, 255, 0)};

\end{tikzpicture}
\end{document}
}
  \end{center}
\end{frame}

% ------------------------------------------------------------------
\begin{frame}{LED 制御 API}
  \begin{table}
    \centering
    \begin{tabular}{lll}
      \hline
      \textbf{関数} & \textbf{説明} & \textbf{引数} \\
      \hline
      \api{disable\_led\_task()} & システム LED 更新を無効化 & --- \\
      \api{enable\_led\_task()}  & システム LED 更新を再有効化 & --- \\
      \api{led\_color(r, g, b)} & LED 色設定 & 各 0--255 \\
      \api{is\_armed()}         & ARM 状態確認 & true / false \\
      \api{battery\_voltage()}  & バッテリー電圧 & 3.0--4.2 V \\
      \hline
    \end{tabular}
  \end{table}

  \vspace{0.5em}

  \begin{block}{WS2812 LED}
    \begin{itemize}
      \item RGB フルカラー（各チャンネル 8bit = 1677万色）
      \item 1本のデータ線で制御（ESP32 RMT ペリフェラル使用）
      \item \texttt{led\_color()} は裏表両方の LED を同じ色に設定
    \end{itemize}
  \end{block}
\end{frame}

% ------------------------------------------------------------------
\begin{frame}{バッテリー電圧と色 / Battery Voltage \& Color}
  \begin{center}
    \begin{tikzpicture}[
        bar/.style={minimum height=8mm, minimum width=4mm, inner sep=0},
      ]
      % Gradient bar
      \foreach \i in {0,...,20} {
        \pgfmathsetmacro{\r}{255*(1-\i/20)}
        \pgfmathsetmacro{\g}{255*\i/20}
        \definecolor{batcol}{RGB}{\r, \g, 0}
        \node[bar, fill=batcol] at (\i*0.4, 0) {};
      }
      % Labels
      \node[below, font=\small] at (0, -0.6) {3.3\,V};
      \node[below, font=\small\bfseries\color{sfred}] at (0, -1.0) {危険};
      \node[below, font=\small] at (4, -0.6) {3.75\,V};
      \node[below, font=\small] at (8, -0.6) {4.2\,V};
      \node[below, font=\small\bfseries\color{sfgreen}] at (8, -1.0) {満充電};
      % Arrow
      \draw[-{Stealth[length=2mm]}, thick] (-0.5, 0) -- (8.5, 0);
    \end{tikzpicture}
  \end{center}

  \vspace{0.5em}

  \begin{itemize}
    \item LiPo バッテリー: 3.3V（空）〜 4.2V（満充電）
    \item \warn{3.3V 以下で使用禁止！}バッテリーが損傷します
  \end{itemize}
\end{frame}

% ------------------------------------------------------------------
\begin{frame}[fragile]{実習: ARM 状態で LED 色を変える / Hands-on}
  \begin{lstlisting}[style=sfcpp, title={user\_code.cpp}]
#include "workshop_api.hpp"
void setup() {
    ws::print("Lesson 3: LED Control");
    ws::disable_led_task();  // Student LED control
}
void loop_400Hz(float dt) {
    if (ws::is_armed()) {
        ws::led_color(0, 255, 0);    // Green
    } else {
        // Battery level as color gradient
        float v = ws::battery_voltage();
        float level = (v - 3.3f) / (4.2f - 3.3f);
        if (level < 0) level = 0;
        if (level > 1) level = 1;
        uint8_t r = 255 * (1.0f - level);
        uint8_t g = 255 * level;
        ws::led_color(r, g, 0);
    }
}
  \end{lstlisting}
\end{frame}

% ------------------------------------------------------------------
\begin{frame}{チェックポイント / Checkpoint}
  \begin{alertblock}{確認事項}
    \begin{itemize}
      \item[$\square$] DISARM 時に赤〜緑のグラデーション表示
      \item[$\square$] ARM すると緑に変わる
      \item[$\square$] バッテリー残量に応じて色が変化する
    \end{itemize}
  \end{alertblock}

  \vspace{1em}

  \begin{block}{次のレッスン}
    Lesson 4: IMU センサー --- 加速度・ジャイロで姿勢を知る
  \end{block}
\end{frame}

\end{document}
